% Manuscript-ready draft material for the DS-MFG QUBO case study.
% Tables use standard LaTeX only. Equations assume amsmath and amssymb, which
% are present in most journal templates. A manuscript template may replace the
% \hline table styling with booktabs if desired.

\section{DS-MFG QUBO Case Study}
\label{sec:ds-mfg-qubo}

This case study evaluates local classical emulations of QAOA and VQE on the
DS-MFG discrete manufacturing QUBO instance using \texttt{QiskitOpt.jl}. The
input data are the exported \texttt{scalars.csv}, \texttt{L\_vector.csv}, and
\texttt{Q\_matrix.csv} files from the DS-MFG archive, with the Gurobi solution
pool used as ground truth for the original integer-programming variables. The
experiments were run in the reproducibility bundle
\texttt{/home/bernalde/repos/QiskitOpt.jl/ds\_mfg\_student\_package\_20260605};
all artifact paths reported below are relative to that bundle.

\subsection{Audit Framing and Evidence Boundary}
\label{subsec:ds-mfg-audit-framing}

The central claim of this case study is practical rather than asymptotic:
current quantum-optimization tooling can be made to sample meaningful
DS-MFG process-design candidates, but only after substantial formulation,
reduction, parameter-search, backend-configuration, and postprocessing work.
The evidence does not support a speedup claim over Gurobi, nor does it support
hardware-performance superiority. Instead, the study audits what practical
value remains after the engineering effort required to use QAOA and VQE is made
explicit.

The guiding question is therefore: for a PSE process-design problem with a known
classical optimum, what do current QAOA/VQE workflows provide beyond a classical
solve, and what workflow burden is required to obtain and interpret those
samples? The answer is organized around three separate evidence tiers:

\begin{table}[t]
\centering
\scriptsize
\caption{Evidence tiers used in the practical quantum-optimization audit.}
\label{tab:ds-mfg-evidence-tiers}
\begin{tabular}{p{0.22\linewidth}p{0.32\linewidth}p{0.34\linewidth}}
\hline
Evidence tier & Role in this case study & Claim supported \\
\hline
Gurobi ground truth & Original 19-flow-variable integer program and solution
pool. & Defines the global optimum, near-optimal flow assignments, and the
reference objective values used for scoring. \\
\hline
Local Aer emulation & \texttt{QiskitOpt.jl}/Qiskit Aer QAOA and VQE runs on
the 36-variable QUBO and the 19-flow quadratic surrogate. & Shows whether
simulated QAOA/VQE sampling distributions place observable mass on useful
repaired flow assignments. \\
\hline
IBM hardware feasibility & Disabled notebook placeholders and fixed-parameter
candidate workflows selected from local evidence. & Identifies hardware-ready
experiments to run later; no IBM hardware performance or superiority result is
claimed here. \\
\hline
\end{tabular}
\end{table}

\subsection{Problem Statement}
\label{subsec:ds-mfg-problem-statement}

The original DS-MFG instance is a binary integer-programming model for a
discrete manufacturing capital-expenditure selection problem. In the Gurobi
export used here the model is named \texttt{capex}; it has 19 binary flow
variables, 20 constraints, 58 nonzeros, and a minimization objective. The Gurobi
solve was run with solution-pool search enabled and returned 36 pool solutions.
We use this pool as classical ground truth over the original 19-variable
decision space.

The computational task in this case study is not merely to recover one bitstring
with low raw QUBO energy. Rather, for each quantum-emulation method $A$, the
task is to generate an empirical distribution
\[
    \widehat{p}_A(x), \qquad x \in \{0,1\}^{36},
\]
then compare its mass against the Gurobi solution pool after projecting onto the
original 19 flow variables and, when needed, repairing the QUBO auxiliary
variables exactly. The reported performance quantities are therefore:
\begin{enumerate}
    \item the best repaired objective sampled by the method;
    \item the number of sampled reads whose flow projection appears in the
    Gurobi pool;
    \item the number of sampled reads whose flow projection is the Gurobi global
    optimum;
    \item the concentration of the distribution on the top repaired flow
    assignments.
\end{enumerate}

The Gurobi global optimum has IP objective $11.7095$ and flow bitstring
\texttt{1001110100111100011} when the flow variables are ordered as
\texttt{f00}, \texttt{f01}, \ldots, \texttt{f18}. The first audit question is
whether local QAOA or VQE emulations place observable probability on this flow.
The second is whether problem-specific postprocessing and reduction make that
probability easier to observe. The third is how much additional engineering
workflow is required relative to solving and interpreting the original Gurobi
model.

\subsection{QUBO Formulation}
\label{subsec:ds-mfg-qubo-formulation}

The QUBO archive provides a scalar file, a linear coefficient vector, and a
quadratic coefficient matrix. Let
\[
    x = (x_1,\ldots,x_{36}) \in \{0,1\}^{36}
\]
denote the complete QUBO bitstring. The objective encoded by the CSV files is
\begin{equation}
    E_{\mathrm{QUBO}}(x)
    =
    a\left(
        \ell^\top x + x^\top Qx + b
    \right),
    \label{eq:ds-mfg-full-qubo}
\end{equation}
with $a=1.0$, $b=601.4761999999997$, $\ell \in \mathbb{R}^{36}$, and
$Q \in \mathbb{R}^{36 \times 36}$. The optimization problem passed to
\texttt{QiskitOpt.jl} is
\begin{equation}
    \min_{x \in \{0,1\}^{36}} E_{\mathrm{QUBO}}(x).
    \label{eq:ds-mfg-full-qubo-problem}
\end{equation}

The 36 QUBO variables split naturally into original variables and reformulation
variables:
\[
    x = (y,z), \qquad
    y \in \{0,1\}^{19}, \quad z \in \{0,1\}^{17}.
\]
Here $y$ is the vector of original DS-MFG flow decisions and $z$ contains the
auxiliary/slack bits introduced by the QUBO reformulation. With the same
partition applied to $\ell$ and $Q$, the QUBO can be written as
\begin{equation}
\begin{aligned}
    E_{\mathrm{QUBO}}(y,z)
    = a\{&
        b + \ell_y^\top y + \ell_z^\top z
        + y^\top Q_{yy}y + y^\top Q_{yz}z \\
        &+ z^\top Q_{zy}y + z^\top Q_{zz}z
    \}.
\end{aligned}
\label{eq:ds-mfg-partitioned-qubo}
\end{equation}
The Gurobi pool is defined only over $y$. Consequently, a sampled bitstring
$(y,z)$ can have a Gurobi-pool flow projection while still carrying a poor
auxiliary assignment $z$ and therefore a misleadingly large raw QUBO energy.

\subsection{Exact Auxiliary Elimination}
\label{subsec:ds-mfg-auxiliary-elimination}

For a fixed flow assignment $y$, define the repaired or reduced flow objective
\begin{equation}
    F(y)
    =
    \min_{z \in \{0,1\}^{17}} E_{\mathrm{QUBO}}(y,z).
    \label{eq:ds-mfg-repaired-objective}
\end{equation}
This objective is the correct bridge between the QUBO formulation and the
original IP solution pool: it asks what QUBO energy is obtained by the best
auxiliary completion of the sampled flow decision. If $y$ appears in the Gurobi
pool, $F(y)$ is also directly comparable to the Gurobi objective value; for the
global optimum above, $F(y^\star)=11.7095$ up to floating-point roundoff.

The reduction is tractable because fixing $y$ removes all flow-variable
uncertainty from Eq.~\eqref{eq:ds-mfg-partitioned-qubo}. The remaining
minimization over $z$ is a binary quadratic problem whose only free variables
are auxiliary bits. In this instance, the auxiliary interaction graph
decomposes into 13 connected components after the flow bits are fixed. Each
component contains either one or two auxiliary variables. Therefore,
\begin{equation}
    F(y)
    =
    c(y) + \sum_{k=1}^{13}
    \min_{z_{C_k} \in \{0,1\}^{|C_k|}}
    \phi_k\!\left(z_{C_k}; y_{N(C_k)}\right),
    \label{eq:ds-mfg-component-repair}
\end{equation}
where $C_k$ is an auxiliary component, $N(C_k)$ is the set of flow variables
that couple to that component, $c(y)$ collects terms independent of auxiliary
variables, and $\phi_k$ is the component-local auxiliary contribution including
the fixed flow-neighbor terms. Since $|C_k| \le 2$ for all components in this
instance, exact repair requires only one- and two-bit local minimizations once a
flow assignment is fixed.

\begin{table}[t]
\centering
\scriptsize
\caption{Auxiliary components used for exact repair after fixing the 19 flow
variables. Flow-neighbor indices use the 1-based indexing stored in the cached
artifact \texttt{auxiliary\_components.csv}.}
\label{tab:ds-mfg-aux-components}
\begin{tabular}{p{0.12\linewidth}p{0.18\linewidth}p{0.30\linewidth}p{0.12\linewidth}p{0.16\linewidth}}
\hline
Component & Auxiliary indices & Flow-neighbor set & Aux. vars. & Flow-neighbor count \\
\hline
1 & 20 & 17;18 & 1 & 2 \\
2 & 21 & 2;3;4 & 1 & 3 \\
3 & 22 & 6;7 & 1 & 2 \\
4 & 23 & 9;10;11 & 1 & 3 \\
5 & 24 & 14;15;16;17 & 1 & 4 \\
6 & 25;26 & 2;7;11 & 2 & 3 \\
7 & 27;28 & 3;7;11 & 2 & 3 \\
8 & 29;30 & 4;7;11 & 2 & 3 \\
9 & 31;32 & 2;3;6 & 2 & 3 \\
10 & 33 & 6;11 & 1 & 2 \\
11 & 34 & 14;18 & 1 & 2 \\
12 & 35 & 15;18 & 1 & 2 \\
13 & 36 & 16;18 & 1 & 2 \\
\hline
\end{tabular}
\end{table}

The notebook evaluates $F(y)$ for all $2^{19}$ flow assignments and verifies
that the exact reduced objective recovers the Gurobi global optimum. This
enumeration also gives a ranked list of the best repaired flow assignments and
is used to score every sampled QAOA and VQE distribution.

Although $F(y)$ is exact, it is not generally quadratic in $y$. The minimization
over auxiliary bits in Eq.~\eqref{eq:ds-mfg-repaired-objective} creates a
piecewise Boolean function of the flow variables. The current
\texttt{QiskitOpt.jl} QAOA and VQE interfaces accept quadratic QUBOs, so the
notebook constructs a quadratic surrogate
\begin{equation}
    \widetilde{F}(y)
    =
    \widetilde{b}
    + \widetilde{\ell}^{\top} y
    + y^{\top} \widetilde{Q} y
    \label{eq:ds-mfg-quadratic-surrogate}
\end{equation}
by least-squares fitting Eq.~\eqref{eq:ds-mfg-quadratic-surrogate} to
$F(y)$ over all $2^{19}$ flow assignments. The fitted surrogate is stored as
\texttt{reduced\_scalars.csv}, \texttt{reduced\_L\_vector.csv}, and
\texttt{reduced\_Q\_matrix.csv}. Its offset is $144.603743749989$ and its
scale is $1.0$.

The surrogate has high global fit quality ($R^2 = 0.9990819$), but the
surrogate optimum is not the repaired global optimum. The fitted surrogate is
minimized by flow \texttt{1001110110011100011}, whose exact repaired objective
is $14.5505$, whereas the true repaired optimum is
\texttt{1001110100111100011} with objective $11.7095$. This is why all
surrogate-based QAOA and VQE samples are evaluated against $F(y)$, not only
against $\widetilde{F}(y)$.

\subsection{Computational Setup}
\label{subsec:ds-mfg-computational-setup}

All QAOA and VQE experiments reported here are local classical simulations, not
IBM hardware executions. The notebook was executed through
\texttt{scripts/nbconvert\_ds\_mfg.sh}, which activates the Julia project
\texttt{DSMFGQiskitOptStudent}. The project pins \texttt{QiskitOpt.jl v0.4.2},
\texttt{QUBODrivers v0.6}, \texttt{QUBOTools v0.13}, \texttt{PythonCall v0.9},
and a Julia compatibility target of \texttt{1.10}. The QAOA and VQE solves use
Qiskit Aer through \texttt{QiskitOpt.jl}. The expensive local simulations use
Aer's matrix-product-state method with single precision; the cached full-QUBO
notebook rows record the fast MPS settings used there, while the reduced QAOA
and VQE revisit scripts explicitly set truncation threshold $10^{-10}$,
maximum MPS bond dimension 64, \texttt{mps\_heuristic} measurement sampling,
and a thread cap of \texttt{min(Sys.CPU\_THREADS, 8)}. Optimizer reads and final
sampling reads were separated with \texttt{QUBODrivers.FinalNumberOfReads()}
when supported by the package.

Because these runs are simulator experiments, timing and hit-rate observations
are not hardware-performance measurements. The reported results should be read
as evidence about workflow feasibility and sampled-distribution quality under
local Aer emulation, not as evidence of quantum speedup or hardware superiority.

The strongest QAOA workflow used an additional offline angle-search step. That
step was run from the local sibling clone of LANL's \texttt{JuliQAOA.jl}. The
script enumerates all $2^{19}$ reduced-surrogate energies for statevector
probability evaluation, normalizes those energies, and then uses
\texttt{find\_angles\_bh} for basin-hopping angle optimization. The resulting
angles were transferred back into \texttt{QiskitOpt.QAOA} for Aer sampling with
\texttt{MaximumIterations() = 0}. Thus, the final QAOA sampling distribution is
still a \texttt{QiskitOpt.jl}/Aer result, but the angles were learned by
\texttt{JuliQAOA.jl}.

\subsection{Distribution Scoring}
\label{subsec:ds-mfg-scoring}

The full sampled state is still retained during scoring. For a sampled state
$x=(y,z)$, the raw QUBO energy is $E_{\mathrm{QUBO}}(y,z)$, the repaired energy
is $F(y)$, and the auxiliary gap is
\[
    \Delta_{\mathrm{aux}}(y,z)
    =
    E_{\mathrm{QUBO}}(y,z) - F(y).
\]
This gap is zero only when the sampled auxiliary bits are already a
minimum-energy completion for the sampled flow projection.

For each sampled state, the notebook reports three related quantities:
\begin{enumerate}
    \item \emph{projected match}: whether the first 19 bits appear in the
    Gurobi pool;
    \item \emph{encoded match}: whether the full 36-bit state matches the
    minimum-energy auxiliary completion for that projected flow;
    \item \emph{repaired objective}: the QUBO objective after fixing the 19
    flow bits and choosing the exact best auxiliary completion.
\end{enumerate}

\subsection{Practical Workflow Burden}
\label{subsec:ds-mfg-workflow-burden}

The practical audit treats the extra workflow required by QAOA and VQE as a
result, not as incidental setup. The Gurobi baseline is solved and interpreted
directly in the original 19 flow variables. By contrast, the quantum-emulation
workflow required importing a 36-variable QUBO, separating flow and auxiliary
bits, repairing auxiliary assignments exactly, enumerating all $2^{19}$ repaired
flow values, fitting a quadratic surrogate because the exact reduced objective
is not generally quadratic, configuring Aer MPS options, separating optimizer
and final sampling reads, running an offline QAOA angle search in a second Julia
package, transferring fixed parameters back into \texttt{QiskitOpt.jl}, and
postprocessing every sampled state against the repaired objective.

This burden is central to the interpretation. The local QAOA and VQE results
show that meaningful samples can be obtained, but they were not obtained by a
drop-in replacement for the classical solver. They required problem-specific
reduction, manual parameter-search choices, backend-specific simulator settings,
and postprocessing logic tied to the original IP variables. The hardware cells
therefore represent feasibility candidates selected from local evidence, not a
completed hardware comparison.

\begin{table}[t]
\centering
\scriptsize
\caption{Experiment campaigns and execution backends. Artifact paths are
relative to the reproducibility bundle.}
\label{tab:ds-mfg-campaigns}
\begin{tabular}{p{0.22\linewidth}p{0.20\linewidth}p{0.28\linewidth}p{0.22\linewidth}}
\hline
Campaign & Problem & Execution settings & Cached evidence \\
\hline
Baseline QAOA/VQE and iteration sweep & Full 36-variable QUBO &
Local \texttt{QiskitOpt.jl}/Aer emulation; small read counts used to verify
the interface and runtime; QAOA iteration limits swept from 3 to 40. &
\texttt{ds\_mfg\_saved\_distributions/};
\texttt{ds\_mfg\_iteration\_sweep/sweep\_summary.csv} \\
\hline
Fixed-angle QAOA, $p=1,\ldots,5$ & Full 36-variable QUBO &
Local Aer MPS; 128 reads, 25 optimizer iterations; Wurtz-Lykov
3-regular MaxCut fixed-angle table used only as a warm start; negative gamma
sign convention. &
\texttt{ds\_mfg\_fixed\_angle\_sweep/fixed\_angle\_qaoa\_summary.csv};
\texttt{ds\_mfg\_aux\_repair\_analysis/repair\_summary.csv} \\
\hline
Perturbed fixed-angle QAOA & Full 36-variable QUBO &
Local Aer MPS; $p=2$ fixed-angle warm start perturbed with
$\sigma \in \{0.05,0.10,0.20\}$; 128 reads, 10 optimizer iterations. &
\texttt{ds\_mfg\_p2\_perturbation\_sweep\_v3/p2\_perturbation\_summary.csv} \\
\hline
QAOA final-sampling separation & Full 36-variable QUBO &
Local Aer MPS; $p=2$; 128 optimizer reads, 512 final reads, 25 optimizer
iterations. &
\texttt{ds\_mfg\_final\_sampling\_sweep\_v2/final\_sampling\_summary.csv} \\
\hline
Auxiliary repair and reduced objective & Exact repaired 19-flow objective &
Deterministic Julia/QUBOTools enumeration of all $2^{19}$ flow assignments;
17 auxiliary variables eliminated by exact component-wise repair. &
\texttt{ds\_mfg\_reduced\_flow\_objective/reduced\_flow\_summary.csv};
\texttt{ds\_mfg\_reduced\_flow\_objective/reduced\_exact\_top\_flows.csv} \\
\hline
Offline QAOA angle search & 19-variable quadratic surrogate &
Local \texttt{JuliQAOA.jl} statevector probability evaluation over all $2^{19}$
surrogate states; basin-hopping angle optimization for $p=1,\ldots,5$;
basin-hopping iterations set to 5 in the cached run; energies z-score
normalized before angle search. &
\texttt{ds\_mfg\_qaoa\_juliqaoa\_angle\_search/juliqaoa\_angle\_summary.csv} \\
\hline
Transferred-angle QAOA & 19-variable quadratic surrogate &
Local \texttt{QiskitOpt.QAOA}/Aer MPS sampling with JuliQAOA angles; 16
optimizer reads, \texttt{MaximumIterations() = 0}; 32768 final reads for
$p=2,\ldots,5$ and 262144 final reads for the best $p=5$ follow-up. &
\texttt{ds\_mfg\_qaoa\_juliqaoa\_transfer/};
\texttt{ds\_mfg\_qaoa\_juliqaoa\_transfer\_highread/} \\
\hline
VQE full-QUBO sweep & Full 36-variable QUBO &
Local \texttt{QiskitOpt.VQE}/Aer emulation with \texttt{EfficientSU2}; five
random initializations; 128 optimizer reads, 8192 final reads, 25 COBYLA
iterations. &
\texttt{ds\_mfg\_vqe\_final\_sampling\_sweep\_v2/vqe\_final\_sampling\_summary.csv} \\
\hline
VQE reduced-surrogate sweeps & 19-variable quadratic surrogate &
Local \texttt{QiskitOpt.VQE}/Aer emulation with \texttt{EfficientSU2}; first
five seeds with 8192 final reads, then 20 seeds with 32768 final reads, then
selected high-read follow-ups at 262144 and 524288 reads. &
\texttt{ds\_mfg\_vqe\_reduced\_flow\_objective\_v2/};
\texttt{ds\_mfg\_vqe\_reduced\_flow\_objective\_v3/};
\texttt{ds\_mfg\_vqe\_reduced\_flow\_objective\_v4/};
\texttt{ds\_mfg\_vqe\_reduced\_flow\_objective\_final/} \\
\hline
Classical sampling baselines & Exact repaired 19-flow objective &
Uniform random sampling over the 19 flow variables and steepest-descent
hill-climbing with random restarts; all sampled flows scored by exact auxiliary
repair. &
\texttt{ds\_mfg\_classical\_baselines/classical\_baseline\_summary.csv};
\texttt{ds\_mfg\_classical\_baselines/classical\_baseline\_distribution.csv} \\
\hline
Hit-rate uncertainty and time to solution & Cached QAOA, VQE, and classical
summary rows &
Two-sided 95\% Wilson hit-rate intervals and empirical 95\% time to solution
for top-50, top-10, and global-optimum events. &
\texttt{ds\_mfg\_hit\_rate\_reports/time\_to\_solution\_report.csv} \\
\hline
\end{tabular}
\end{table}

\subsection{QAOA Results}
\label{subsec:ds-mfg-qaoa-results}

The initial full-QUBO QAOA experiments did not concentrate probability mass on
encoded Gurobi-pool solutions. Fixed-angle QAOA with $p=2$ produced one
projected Gurobi-pool read in 128 samples, but no encoded pool read. After
exact auxiliary repair, that projected flow had objective $18.0745$, well above
the global optimum. Perturbing those $p=2$ fixed angles did not improve the
result. Separating optimizer reads from final sampling on the full QUBO improved
the best repaired pool objective to $14.5505$ and produced three projected pool
reads in 512 samples, but still did not sample the global optimum.

The strongest QAOA result was obtained by optimizing angles offline on the
reduced surrogate with \texttt{JuliQAOA.jl} and then transferring those angles
into \texttt{QiskitOpt.QAOA}. For $p=5$, the offline statevector calculation
predicted a probability of $0.1506823$ on the top 50 repaired flows, $0.0340277$
on the top 10 repaired flows, and $0.00263256$ on the exact global optimum. The
high-read Aer follow-up sampled the global optimum 664 times in 262144 reads,
corresponding to an empirical global-hit rate of $0.002533$.

\begin{table}[t]
\centering
\scriptsize
\caption{QAOA experiments scored by exact auxiliary repair. ``Global'' means
the Gurobi global optimum \texttt{1001110100111100011}.}
\label{tab:ds-mfg-qaoa-results}
\begin{tabular}{p{0.25\linewidth}p{0.19\linewidth}p{0.22\linewidth}p{0.25\linewidth}}
\hline
Experiment & Settings & Hit counts & Best repaired result \\
\hline
Fixed-angle QAOA, best full-QUBO case & $p=2$; 128 reads; 25 iterations; full
36-variable QUBO & 1 projected pool read; 0 encoded pool reads; 0 global reads &
Objective $18.0745$; flow \texttt{1001110101011100011}; artifact
\texttt{repair\_summary.csv} \\
\hline
Perturbed $p=2$ fixed-angle QAOA & $\sigma=0.05,0.10,0.20$; 128 reads; 10
iterations; full 36-variable QUBO & 0 projected pool reads and 0 global reads
for all three perturbations & Best repaired objective $86.8695$ at
$\sigma=0.10$; artifact \texttt{p2\_perturbation\_summary.csv} \\
\hline
Full-QUBO QAOA with separated final sampling & $p=2$; 128 optimizer reads; 512
final reads; 25 iterations & 3 projected pool reads; 0 encoded pool reads; 0
global reads & Objective $14.5505$; flow
\texttt{1001110110011100011}; artifact
\texttt{final\_sampling\_summary.csv} \\
\hline
Reduced-surrogate QAOA before offline angle transfer & $p=2$; 128 optimizer
reads; 512 final reads; 25 iterations & 0 pool reads and 0 global reads &
Best sampled repaired objective $102.7385$; artifact
\texttt{reduced\_flow\_summary.csv} \\
\hline
Transferred-angle QAOA, sweep & $p=2,\ldots,5$; 32768 final reads per depth;
\texttt{MaximumIterations() = 0} & Global reads increased from 16 at $p=2$ to
96 at $p=5$ & Every depth sampled the global optimum at least once; artifact
\texttt{qaoa\_juliqaoa\_transfer\_summary.csv} \\
\hline
Transferred-angle QAOA, high-read follow-up & $p=5$; 16 optimizer reads;
262144 final reads; \texttt{MaximumIterations() = 0} & 39661 top-50 reads;
8964 top-10 reads; 664 global reads & Objective $11.7095$; flow
\texttt{1001110100111100011}; artifact
\texttt{qaoa\_juliqaoa\_transfer\_highread/qaoa\_juliqaoa\_transfer\_summary.csv} \\
\hline
\end{tabular}
\end{table}

The $p=5$ transferred QAOA parameters are stored in beta-then-gamma order after
undoing the JuliQAOA energy normalization:
\begin{verbatim}
2.33962976942;2.63487051398;2.8837299086;2.94861783352;3.01149134351;
0.00156283712269;0.00374327754686;0.00643793480657;0.0100189310157;
0.0141285087419
\end{verbatim}
These are the fixed parameters used by the QAOA hardware placeholder in the
notebook.

\subsection{VQE Results}
\label{subsec:ds-mfg-vqe-results}

The initial VQE experiments were less concentrated than the best QAOA workflow.
On the full 36-variable QUBO, a five-seed \texttt{EfficientSU2} sweep with
8192 final reads per seed produced one projected pool read and no encoded or
global read. The best repaired full-QUBO VQE objective was $14.6515$.

The reduced-surrogate VQE runs performed better. A five-seed sweep with 8192
final reads found a near-optimal pool flow with repaired objective $11.8105$,
but no global read. A 20-seed sweep with 32768 final reads per seed then
sampled the global optimum three times in 655360 total reads. Subsequent
high-read follow-ups on seeds 74001, 74007, and 74018 increased the total
number of observed global reads to 19 at 262144 reads per seed, and to 28 at
524288 reads per seed. The final follow-up identified seed 74018 as the best
VQE candidate, with 20 global reads, 152 top-10 reads, and 391 top-50 reads in
524288 final samples.

\begin{table}[t]
\centering
\scriptsize
\caption{VQE experiments scored by exact auxiliary repair. All VQE runs used
\texttt{EfficientSU2} and local \texttt{QiskitOpt.VQE}/Aer emulation.}
\label{tab:ds-mfg-vqe-results}
\begin{tabular}{p{0.26\linewidth}p{0.20\linewidth}p{0.22\linewidth}p{0.23\linewidth}}
\hline
Experiment & Settings & Hit counts & Best repaired result \\
\hline
Full-QUBO VQE, five seeds & Seeds 73001--73005; 128 optimizer reads; 8192
final reads per seed; 25 iterations & 1 projected pool read; 0 encoded pool
reads; 0 global reads & Objective $14.6515$ from seed 73005; artifact
\texttt{vqe\_final\_sampling\_summary.csv} \\
\hline
Reduced-surrogate VQE, five seeds & Seeds 73001--73005; 128 optimizer reads;
8192 final reads per seed; 25 iterations & 2 pool reads total; 0 global reads &
Objective $11.8105$ from seed 73003; flow
\texttt{1001110100111010011}; artifact
\texttt{vqe\_reduced\_flow\_summary.csv} \\
\hline
Reduced-surrogate VQE, 20 seeds & Seeds 74001--74020; 32768 final reads per
seed; 25 iterations & 3 global reads in 655360 total reads; seed 74018 had the
largest top-50 count but no global read at this read count & Objective
$11.7095$ observed at seeds 74001 and 74007; artifact
\texttt{vqe\_reduced\_top50\_sampling\_summary.csv} \\
\hline
Reduced-surrogate VQE, high-read selected seeds & Seeds 74001, 74007, 74018;
262144 final reads per seed & 19 global reads in 786432 total reads; seed
74018 produced 13 global reads & Objective $11.7095$; artifact
\texttt{ds\_mfg\_vqe\_reduced\_flow\_objective\_v4/} \\
\hline
Reduced-surrogate VQE, final follow-up & Seeds 74018, 74007, 74001; 524288
final reads per seed & 28 global reads in 1572864 total reads; seed 74018
produced 20 global reads, 152 top-10 reads, and 391 top-50 reads & Objective
$11.7095$; artifact
\texttt{ds\_mfg\_vqe\_reduced\_flow\_objective\_final/} \\
\hline
\end{tabular}
\end{table}

\subsection{Classical Sampling Baselines}
\label{subsec:ds-mfg-classical-baselines}

To interpret QAOA and VQE hit rates against non-quantum alternatives, two
classical baselines were run directly on the repaired 19-flow objective
$F(y)$. The first baseline samples flow assignments uniformly at random and
then applies exact auxiliary repair. The second baseline uses steepest-descent
single-bit-flip hill climbing with random restarts; each evaluated flow counts
against the same sample/evaluation budget used in the table. Both baselines use
fixed seeds and the same top-50, top-10, and global-optimum definitions used for
the QAOA and VQE summaries. The cached outputs include a compressed aggregate
distribution over every unique sampled flow, a compact retained-flow file for
sampled top-50 and best-sampled flows, and a summary that records counts over
the full budget.

\begin{table}[t]
\centering
\scriptsize
\caption{Classical baseline context for the strongest local QAOA and VQE
sampling results. ``Samples'' are reads for QAOA/VQE and repaired-flow
objective evaluations for the classical baselines.}
\label{tab:ds-mfg-classical-baseline-context}
\begin{tabular}{p{0.24\linewidth}p{0.19\linewidth}p{0.23\linewidth}p{0.24\linewidth}}
\hline
Method & Samples / evaluations & Top-hit counts & Best repaired result \\
\hline
Uniform random, seed 81001 & 262144 & 24 top-50; 5 top-10; 0 global &
Objective $11.8105$; rank 2; flow \texttt{1001110100111010011} \\
\hline
Uniform random, seed 81002 & 524288 & 49 top-50; 11 top-10; 0 global &
Objective $12.2415$; rank 3; flow \texttt{1001110100111001011} \\
\hline
Hill-climb restarts, seed 82001 & 262144 evaluations; 2055 restarts &
883 top-50; 177 top-10; 19 global &
Objective $11.7095$; rank 1; flow \texttt{1001110100111100011} \\
\hline
Transferred-angle QAOA, $p=5$ & 262144 local Aer reads &
39661 top-50; 8964 top-10; 664 global &
Objective $11.7095$; rank 1; flow \texttt{1001110100111100011} \\
\hline
Reduced-surrogate VQE, seed 74018 & 524288 local Aer reads &
391 top-50; 152 top-10; 20 global &
Objective $11.7095$; rank 1; flow \texttt{1001110100111100011} \\
\hline
\end{tabular}
\end{table}

The cached QAOA, reduced-surrogate VQE, and classical-baseline summary CSVs
also report empirical hit rates for the top-50, top-10, and global-optimum
events. Each rate is paired with a two-sided 95\% Wilson score confidence
interval. The same rows report empirical 95\% time to solution,
\texttt{tts95\_sec}, computed as the recorded local solve or wall time per
sample multiplied by
$\lceil \log(1 - 0.95) / \log(1 - \hat{p}) \rceil$, where $\hat{p}$ is the
empirical hit rate. Rows with zero observed hits are reported with infinite
empirical time to solution rather than an extrapolated finite value. QAOA and
VQE timing uses local Aer \texttt{solve\_time\_sec}; the classical baselines
use script \texttt{wall\_time\_sec}. The consolidated artifact is
\texttt{ds\_mfg\_hit\_rate\_reports/time\_to\_solution\_report.csv}.

\begin{table}[t]
\centering
\scriptsize
\caption{Representative global-hit uncertainty and empirical 95\% time to
solution from cached local runs.}
\label{tab:ds-mfg-global-hit-uncertainty}
\begin{tabular}{p{0.30\linewidth}p{0.17\linewidth}p{0.30\linewidth}p{0.13\linewidth}}
\hline
Run & Global hits & Global hit rate, 95\% Wilson interval & \texttt{tts95\_sec} \\
\hline
Transferred $p=5$ QAOA, high-read & 664 / 262144 &
$0.002533$ [$0.002348$, $0.002733$] & $0.243$ \\
\hline
Reduced-surrogate VQE, seed 74018 & 20 / 524288 &
$0.00003815$ [$0.00002470$, $0.00005892$] & $4.505$ \\
\hline
Hill-climb restarts, seed 82001 & 19 / 262144 &
$0.00007248$ [$0.00004640$, $0.00011321$] & $0.147$ \\
\hline
Uniform random, seed 81001 & 0 / 262144 &
$0$ [$0$, $0.00001465$] & $\infty$ \\
\hline
\end{tabular}
\end{table}

Uniform random sampling provides a useful lower baseline for concentration. At
262144 samples it found 24 top-50 flows, 5 top-10 flows, and no global optimum
read; at 524288 samples it found 49 top-50 flows, 11 top-10 flows, and again no
global optimum read. Both reduced-surrogate QAOA and the selected VQE follow-up
therefore show much stronger concentration on low repaired-objective flows than
uniform random sampling at comparable sample counts.

The hill-climbing baseline is a stronger classical comparator because it uses
the repaired objective adaptively. With the same 262144 evaluation budget used
for the QAOA high-read follow-up, it found the global optimum 19 times and
reached 883 top-50 evaluations. This confirms that simple classical search can
also exploit $F(y)$ once exact repair is available. However, the transferred
$p=5$ QAOA distribution was still more concentrated in this cached comparison,
with 664 global reads and 39661 top-50 reads at the same nominal count. The
comparison does not imply a runtime advantage for QAOA; it shows that sampled
distribution concentration should be interpreted relative to simple classical
sampling baselines, not only relative to Gurobi.

\subsection{Interpretation}
\label{subsec:ds-mfg-interpretation}

The dominant methodological point is that a process-design quantum workflow
must preserve a clear bridge back to the original optimization model. The
original IP variables and the QUBO auxiliary variables must be separated in the
analysis. A sample can be useful at the IP-variable level even when its
auxiliary bits give a poor raw QUBO energy. Exact auxiliary repair therefore
provides the correct comparison to the Gurobi solution pool and the original IP
objective.

The dominant algorithmic point is that transferred-angle QAOA was the strongest
workflow tested. Direct local optimization of the full 36-variable QUBO was weak
for both QAOA and VQE. Reducing the problem to the 19 flow variables and
repairing auxiliary variables exactly made the problem interpretable, and
offline statevector QAOA angle search on the reduced surrogate produced a
distribution with visible mass on the global optimum. VQE also reached the
global optimum on the reduced surrogate, but at a much lower empirical rate:
the best final VQE follow-up saw 20 global reads in 524288 samples for seed
74018, whereas the transferred $p=5$ QAOA follow-up saw 664 global reads in
262144 samples.

The classical baselines add sampling context to this conclusion. Uniform random
sampling almost never reaches the global optimum at the budgets used here, so
both the QAOA transfer result and the selected VQE follow-up are meaningfully
more concentrated than random repaired-flow sampling. A simple hill-climbing
baseline can also reach the global optimum using the repaired objective, which
underscores that exact repair and objective access are powerful classical tools
in their own right. The quantum-emulation results should therefore be framed as
sampled-distribution evidence after substantial preprocessing, not as evidence
that quantum optimization dominates simple classical search.

The practical conclusion is more conservative than the algorithmic comparison:
the DS-MFG instance can be reformulated and sampled with current QAOA/VQE
software, and local simulations can recover the Gurobi optimum after reduction
and repair, but the workflow burden is substantial. The evidence supports an
audit claim about feasibility, interpretation, and engineering effort; it does
not support replacing Gurobi for this process-design instance.

\subsection{Hardware Run Plan}
\label{subsec:ds-mfg-hardware-plan}

The notebook includes placeholder cells for hardware execution, but no hardware
results are included in this bundle. The QAOA hardware candidate is ready as a
fixed-parameter sampling experiment. The VQE hardware candidate is less mature
because the cache identifies the best random seed and sampled distribution, but
does not yet persist the optimized \texttt{EfficientSU2} parameter vector.
These placeholders are included to make the next experimental step explicit;
they are not evidence of IBM hardware performance.

\begin{table}[t]
\centering
\scriptsize
\caption{Hardware placeholders recorded in the notebook.}
\label{tab:ds-mfg-hardware-placeholders}
\begin{tabular}{p{0.20\linewidth}p{0.30\linewidth}p{0.38\linewidth}}
\hline
Method & Local evidence supporting the candidate & Placeholder hardware action \\
\hline
QAOA & Reduced 19-flow surrogate; $p=5$; JuliQAOA-learned angles transferred
to \texttt{QiskitOpt.QAOA}; 664 global reads in 262144 local Aer samples. &
Submit a fixed-angle $p=5$ QAOA sampling run with
\texttt{MaximumIterations() = 0} and the beta-then-gamma parameter vector
listed in Section~\ref{subsec:ds-mfg-qaoa-results}. \\
\hline
VQE & Reduced 19-flow surrogate; \texttt{EfficientSU2}; best final follow-up
seed 74018; 20 global reads in 524288 local Aer samples. & Rerun VQE from
seed 74018 unless the scripts are extended to persist and reload the optimized
ansatz parameter vector. For a hardware experiment, persisting that vector
should be done before submission so optimization and hardware sampling can be
separated. \\
\hline
\end{tabular}
\end{table}

\subsection{Limitations}
\label{subsec:ds-mfg-limitations}

The local results are classical Aer simulations and should not be interpreted as
hardware performance. The reduced objective is exact after auxiliary repair,
but the QAOA and VQE experiments use a quadratic surrogate fit to that exact
objective. Although the surrogate has high global $R^2$, its minimizer is not
the true repaired global optimum. Future work should test either an interface
that optimizes the exact reduced objective directly or a surrogate fit weighted
toward the low-energy region of the repaired objective. The VQE workflow should
also persist optimized ansatz parameters so that future hardware studies can
run fixed-parameter sampling rather than rerunning the classical optimization
loop on hardware.
